
          %%%%% ~~~~~~~~~~~~~~~~~~~~ %%%%%

\chapter{Blackbody radiation}
\setcounter{theorem}{0}
\setcounter{equation}{0}

%\cite{vanDerVaart1996}
%\cite{Kosorok2008}

%\renewcommand{\theenumi}{\alph{enumi}}
%\renewcommand{\labelenumi}{\textnormal{(\theenumi)}$\;\;$}
\renewcommand{\theenumi}{\roman{enumi}}
\renewcommand{\labelenumi}{\textnormal{(\theenumi)}$\;\;$}

          %%%%% ~~~~~~~~~~~~~~~~~~~~ %%%%%

\section{Universality of blackbody emission spectrum -- Kirchhoff's law of thermal radiation}

          %%%%% ~~~~~~~~~~~~~~~~~~~~ %%%%%

\vskip 0.5cm
\begin{theorem}[Universality of the Blackbody Emission Spectrum]
\mbox{}
\vskip 0.1cm
\noindent
Suppose:
\begin{itemize}
\item
    $\mathcal{C}$\, is the set of all possible macroscopic blackbody configurations
    (encompassing variations in shape, volume, and interior wall material).
\item
    $T \in \Re^{+}$ is the absolute temperature, and
    \,$\nu \in \Re^{+}$ is the frequency of electromagnetic radiation.
\item
    \,$u: \mathbb{R}^+ \times \mathbb{R}^+ \times \mathcal{C} \xrightarrow{{\color{white}....}} \mathbb{R}^+$\,
    is the blackbody emission spectral energy density function,
    i.e., \,$u(\nu, T, C)$\, is the blackbody emission energy density per unit frequency
    at frequency \,$\nu \in \Re^{+}$,\, temperature \,$T \in \Re^{+}$,\,
    for macroscopic blackbody configuration \,$C \in \mathcal{C}$.\,
\end{itemize}
Assume:
\begin{itemize}
\item
    \textbf{Axiom 1 (Continuity):}\;
    For any fixed $T$ and $C$, the spectral energy density $u(\nu, T, C)$ is a continuous function of the frequency $\nu$.
\item
    \textbf{Axiom 2 (Strict Monotonicity of Flux):}\;
    The spectral power per unit area $\Phi_\nu$ emitted from a small aperture in cavity $C$ is a strictly monotonically increasing function of its internal spectral energy density. That is, $\Phi_\nu = f(u(\nu, T, C))$, where $f: \mathbb{R}^+ \to \mathbb{R}^+$ satisfies $f(a) > f(b) \iff a > b$. (In classical electrodynamics, $f(u) = \frac{c}{4}u$, but strict monotonicity is mathematically sufficient).
\item
    \textbf{Axiom 3 (Second Law of Thermodynamics):}\;
    If two cavities $C_1, C_2 \in \mathcal{C}$ are maintained at the exact same temperature $T$ and are allowed to exchange energy radiatively, the net steady-state heat flux between them must be exactly zero.
\item
    \textbf{Axiom 4 (Ideal Optical Filtration):}\;
    For any interval $[\nu_1, \nu_2] \subset \mathbb{R}^+$, it is possible to construct a filter that transmits $100\%$ of incident radiation within $[\nu_1, \nu_2]$ and reflects $100\%$ of radiation outside this interval.
\end{itemize}
Then, the following statement is true:
For any two arbitrary macroscope blackbody configurations
\,$C_1, C_2 \,\in\, \mathcal{C}$\,
at the same temperature \,$T > 0$,\,
the spectral energy density is identical for all frequencies, i.e.,
\begin{equation*}
u(\nu, T, C_1)
\;\;=\;\;
    u(\nu, T, C_2),
\quad
\textnormal{for each \,$\nu > 0$,\, each \,$T > 0$,\, and each \,$C_{1}, C_{2} \in \mathcal{C}$}
\end{equation*}
Consequently, \,$u$\, is independent of \,$C$\, and is a universal function \,$u(\nu, T)$.\,
\end{theorem}
\proof
Assume, for the sake of contradiction, that there exists a specific temperature $T > 0$, a frequency $\nu_0 > 0$, and two configurations $C_1, C_2 \in \mathcal{C}$ such that their spectral energy densities differ.

\vskip 0.2cm
\noindent
Without loss of generality, assume:$$u(\nu_0, T, C_1) > u(\nu_0, T, C_2)$$Because $u$ is continuous with respect to $\nu$ (Axiom 1), there must exist a finite frequency neighborhood $[\nu_1, \nu_2]$ containing $\nu_0$ (where $\nu_1 < \nu_2$) such that for all $\nu \in [\nu_1, \nu_2]$, the inequality holds:$$u(\nu, T, C_1) > u(\nu, T, C_2)$$Let the two cavities $C_1$ and $C_2$ be placed facing each other, connected only by a small aperture of area $A$.Between the cavities, insert an ideal filter (Axiom 4) that exclusively transmits radiation in the frequency range $[\nu_1, \nu_2]$, reflecting all other frequencies back into their respective cavities.The total power $P_{1 \to 2}$ radiated from cavity $C_1$ to $C_2$ through the filter is the integral of the emitted flux over the transmitted frequencies:$$P_{1 \to 2} = A \int_{\nu_1}^{\nu_2} f(u(\nu, T, C_1)) \, d\nu$$Similarly, the total power radiated from cavity $C_2$ to $C_1$ is:$$P_{2 \to 1} = A \int_{\nu_1}^{\nu_2} f(u(\nu, T, C_2)) \, d\nu$$By Axiom 2, because $u(\nu, T, C_1) > u(\nu, T, C_2)$ everywhere on the interval $[\nu_1, \nu_2]$, it strictly follows that $f(u(\nu, T, C_1)) > f(u(\nu, T, C_2))$ on this interval.Since the integral of a strictly positive difference over a non-zero interval $[\nu_1, \nu_2]$ is strictly positive, we have:$$P_{1 \to 2} > P_{2 \to 1}$$This implies there is a net transfer of heat energy $P_{\text{net}} = P_{1 \to 2} - P_{2 \to 1} > 0$ flowing continuously from cavity $C_1$ to cavity $C_2$.However, both $C_1$ and $C_2$ are precisely at the same temperature $T$. A spontaneous, continuous net flow of heat between two bodies at identical temperatures constitutes a violation of the Second Law of Thermodynamics (Axiom 3).Therefore, our initial assumption must be false.We conclude that no such $\nu_0$ exists, and $u(\nu, T, C_1) = u(\nu, T, C_2)$ for all $\nu, T > 0$ and all $C_1, C_2 \in \mathcal{C}$.
\qed

          %%%%% ~~~~~~~~~~~~~~~~~~~~ %%%%%

\section{The ultraviolet catastrophe -- the Rayleigh–Jeans law}

          %%%%% ~~~~~~~~~~~~~~~~~~~~ %%%%%

\section{Explanation of the universal blackbody emission spectrum via discretization by Planck}

          %%%%% ~~~~~~~~~~~~~~~~~~~~ %%%%%

